\subsection{Jet manifolds}

12.3.1. Let \(c = (\mathrm{U}, \varphi, \mathrm{E})\) and \(c' = (\mathrm{V}, \psi, \mathrm{F})\) be charts of X and Y, respectively. The maps \(\varphi\) and \(\psi\) define, by transport of structure, a bijection \(\pi\) of \(\mathbf{J}^k (\mathbf{U},\mathbf{V})\) onto

\[
\mathrm{J} ^ {k} (\varphi (\mathrm{U}), \psi (\mathrm{V})) = \varphi (\mathrm{U}) \times \psi (\mathrm{V}) \times \mathrm{Q} ^ {k} (\mathrm{E}, \mathrm{F}),
\]

cf. 12.2.2. If one sets \(\mathbf{W} = \mathbf{J}^{k}(\mathbf{U},\mathbf{V})\) and \(\mathbf{G} = \mathbf{E}\times \mathbf{F}\times \mathbf{Q}^{k}(\mathbf{E},\mathbf{F})\), the triple \((\mathbf{W},\pi ,\mathbf{G})\) is a chart of \(\mathbf{J}^k (\mathbf{X},\mathbf{Y})\). The charts thus obtained form a \(\mathbf{C}^{r - k}\)-atlas and this endows \(\mathbf{J}^k (\mathbf{X},\mathbf{Y})\) with a K-manifold structure of class \(\mathbf{C}^{r - k}\).\footnote{When \(k=r\) (which is possible only if \(K=\mathbf{R}\)), \(\mathbf{J}^{k}(\mathbf{X},\mathbf{Y})\) is endowed with a topological manifold structure (cf. Notations and Conventions), and the morphisms and fibrations considered below are to be understood in the purely topological sense (cf. the footnote to § 6).}

If \((x, y) \in X \times Y\), the sets \(J_{x}^{k}(X, Y)\), \(J^{k}(X, Y)_{y}\), and \(J_{x}^{k}(X, Y)_{y}\) are closed submanifolds of \(J^{k}(X, Y)\).

12.3.2. If X and Y are pure (resp. finite-dimensional, resp. separated, resp. connected), then so is \(J^{k}(X, Y)\).

If U (resp. V) is open in X (resp. in Y), \(J^{k}(U, V)\) is an open submanifold of \(J^{k}(X, Y)\), cf. 12.1.3.

12.3.3. The maps \(s: \mathbf{J}^{k}(\mathbf{X}, \mathbf{Y}) \to \mathbf{X}, b: \mathbf{J}^{k}(\mathbf{X}, \mathbf{Y}) \to \mathbf{Y}\) and \((s, b): \mathbf{J}^{k}(\mathbf{X}, \mathbf{Y}) \to \mathbf{X} \times \mathbf{Y}\) are fibrations (6.1.1) of class \(\mathbf{C}^{r - k}\) . When \(k = 0\) , \((s, b)\) is an isomorphism.

12.3.4. Denote by \(T_{X}\) and \(T_{Y}\) the vector bundles \(\mathrm{pr}_{1}^{*}\mathrm{T}(\mathrm{X})\) and \(\mathrm{pr}_{2}^{*}\mathrm{T}(\mathrm{Y})\) on \(X \times Y\), and let \(\mathcal{L}(T_{X}; T_{Y})\) be the bundle of homomorphisms of \(T_{X}\) to \(T_{Y}\) (7.7.3); if \((x, y) \in X \times Y\), one has

\[
\mathcal {L} (\mathrm{T} _ {\mathrm{X}}; \mathrm{T} _ {\mathrm{Y}}) _ {(x, y)} = \mathcal {L} (\mathrm{T} _ {x} (\mathrm{X}); \mathrm{T} _ {y} (\mathrm{Y})).
\]

When \(k = 1\), the map

\[
\mathbf {T} \colon \mathbf {J} ^ {1} (\mathrm{X}, \mathrm{Y}) \rightarrow \mathcal {L} (\mathrm{T} _ {\mathrm{X}}; \mathrm{T} _ {\mathrm{Y}})
\]

thus defined is an isomorphism of manifolds of class \(C^{r-1}\).

For X = K, this isomorphism identifies \(J_{0}^{1}(K, Y)\) with T(Y); for Y = K, it identifies \(J^{1}(X, K)_{0}\) with the dual \(T'(X)\) of T(X).

12.3.5. Let \(k'\) be an integer such that \(0 \leqslant k' \leqslant k\) . The map

\[
r ^ {k, k ^ {\prime}}: \mathrm{J} ^ {k} (\mathrm{X}, \mathrm{Y}) \rightarrow \mathrm{J} ^ {k ^ {\prime}} (\mathrm{X}, \mathrm{Y})\tag{cf. 12.1.5}
\]

is a fibration of class \(C^{r-k}\) .

12.3.6. Let Z be a manifold of class \(\mathbf{C}^r\) , and let T be the set of pairs

\[
(j, j ^ {\prime}) \in \mathrm{J} ^ {k} (\mathrm{X}, \mathrm{Y}) \times \mathrm{J} ^ {k} (\mathrm{Y}, \mathrm{Z})
\]

such that \(b(j) = s(j')\). Then T is a submanifold of \(J^k(X, Y) \times J^k(Y, Z)\), and the map \((j, j') \mapsto j' \circ j\) is a morphism of class \(C^{r-k}\) from T to \(J^k(X, Z)\).

12.3.7. If \(f: \mathbf{X} \to \mathbf{Y}\) is a morphism of class \(\mathbf{C}^r\), the map \(x \mapsto j_x^k(f)\) is a morphism \(j^k(f)\) of class \(\mathbf{C}^{r-k}\) from \(\mathbf{X}\) to \(\mathbf{J}^k(\mathbf{X}, \mathbf{Y})\).

12.3.8. Let \(k'\) and \(k''\) be positive integers with sum \(k\). Let \(x \in \mathbf{X}\), let \(\mathbf{U}\) be an open neighborhood of \(x\), and let \(f: \mathbf{U} \to \mathbf{Y}\) be a morphism of class \(\mathbf{C}^{r}\). The map

\[
x \mapsto j _ {x} ^ {k ^ {\prime}} (f) \colon \mathrm{U} \to \mathrm{J} ^ {k ^ {\prime}} (\mathrm{X}, \mathrm{Y})
\]

is of class \(\mathbf{C}^{r - k'}\), and its jet of order \(k''\) at \(x\) depends only on \(j_x^k (f)\). We thus obtain a canonical map

\[
\alpha \colon \mathrm{J} ^ {k} (\mathrm{X}, \mathrm{Y}) \rightarrow \mathrm{J} ^ {k ^ {\prime \prime}} (\mathrm{X}, \mathrm{J} ^ {k ^ {\prime}} (\mathrm{X}, \mathrm{Y}))
\]

which is of class \(C^{r-k}\); if K is of characteristic zero, it is an embedding.

12.3.9. Let \(\mathbf{X}'\), \(\mathbf{Y}'\) be varieties of class \(\mathbf{C}^r\), let \(f: \mathbf{X} \to \mathbf{X}'\) and \(g: \mathbf{Y}' \to \mathbf{Y}\) be morphisms, and let \((x, y') \in \mathbf{X} \times \mathbf{Y}'\), \(x' = f(x)\), \(y = g(y')\). If \(u \in \mathrm{J}_x^k(\mathbf{X}', \mathbf{Y}')_{y'}\), set

\[
\mathrm{J} _ {x ^ {\prime}} ^ {k} (f, g) _ {y} (u) = j _ {y ^ {\prime}} ^ {k} (g) \circ u \circ j _ {x} ^ {k} (f).
\]

One thus obtains a map

\[
\mathrm{J} ^ {k} (f, g) \colon \mathrm{J} ^ {k} (\mathrm{X} ^ {\prime}, \mathrm{Y} ^ {\prime}) \times_ {\mathrm{X} ^ {\prime}} \mathrm{X} \rightarrow \mathrm{J} ^ {k} (\mathrm{X}, \mathrm{Y})
\]

which is of class \(C^{r-k}\).

12.3.10. Let A be a compact subset of X. Let \(\mathcal{C}^{k}(X;Y)\) be the set of maps of class \(C^{k}\) from X to Y. For every \(f\in\mathcal{C}^{k}(X;Y)\), the map \(j^{k}(f)|A\) belongs to the set \(\mathcal{C}(A;J^{k}(X,Y))\) of continuous maps from A to \(J^{k}(X,Y)\). We have thus defined a map \(\lambda\) from \(\mathcal{C}^{k}(X;Y)\) to \(\mathcal{C}(A;J^{k}(X,Y))\). The inverse image under \(\lambda\) of the topology of compact convergence on \(\mathcal{C}(A;J^{k}(X,Y))\) (TG, X, § 3, def. 1) is a topology on \(\mathcal{C}^{k}(X;Y)\); it is called the topology of \(C^{k}\)-uniform convergence on A.

